% concatenated from natmatdioph.tex
\documentclass{article}
\usepackage{a4wide}
\usepackage[utf8]{inputenc}
\usepackage[T1]{fontenc}
\usepackage{amsmath, amssymb, amsthm}
\usepackage{graphicx}
\usepackage{hyperref}
\usepackage[capitalize]{cleveref} % must be loaded after hyperref
\usepackage{tabularx}
\usepackage{url}
\usepackage{xcolor}

% Theorem-like environments

\theoremstyle{plain}
  \newtheorem{theorem}{Theorem}[section]
  \newtheorem{corollary}{Corollary}[section]
  \newtheorem{lemma}{Lemma}[section]
  \newtheorem{proposition}{Proposition}[section]
  \newtheorem*{theorem*}{Theorem}

\theoremstyle{definition}
  \newtheorem{definition}{Definition}[section]
  \newtheorem{remark}{Remark}[section]

\crefname{theorem}{Theorem}{Theorems}
\crefname{theorem*}{Theorem}{Theorems}
\crefname{lemma}{Lemma}{Lemmas}
\crefname{proposition}{Proposition}{Propositions}
\crefname{definition}{Definition}{Definitions}
\crefname{remark}{Remark}{Remarks}


\title{On polynomial systems of equations in square matrices filled with natural numbers}
\author{Mihai Prunescu}
\date{}



\begin{document}

\maketitle

\begin{abstract}
The positive existential theories of the sets $M_n(\mathbb N)$ without parameters build an inclusion lattice isomorhic with the lattice of divisibility. All these sets are algorithmically undecidable. In further sections some easier observations are made, like the undecidability of diophantine equations with coefficients in $\mathbb M_n(\mathbb Z)$.

Mathematics Subject Classification:  11U55, 11D72, 11C20, 15B48
\end{abstract}

%%%%%%%%%%%%%%%%%%%%%%%%%%%%%%
\section{Introduction}\label{sect:introduction} 

In a list of problems moderated by B. Poonen and T. Scanlon and written down by J. Demeyer, question number 7 (M. Davis) \cite{Demeyer} asks about the existence of decision procedures for non-commutative polynomial equations over structures like the field of quaternions, its ring of integers, or matrix rings like $M_n(\mathbb Q)$ or $M_n(\mathbb Z)$. This paper deals mainly with the set $M_n(\mathbb N)$. This set does not build a ring, but is somewhat analogous with the set of natural numbers $\mathbb N$. In this paper $0$ is considered a natural number. 

The ring of {\bf non-commutative polynomials} $\mathbb Z[X_1,X_2, \dots, X_k]$ is built exactly like the usual ring of polynomials, but contains more monomials. So, we will explicitly use monomials like $AYA$, which is different from $A^2Y$ and also from $YA^2$. Also, we will explicitly use the {\bf commutator} polynomial $XY-YX$ which is different from the polynomial zero.  

The Theorem of Matiyasevich \cite{Matiyasevich} says that given $P \in \mathbb Z[X_1,X_2, \dots, X_k]$, it is undecidable whether 
$$\exists \, X_1, X_2, \dots, X_k \in \mathbb N\,\,\,\,P(X_1,X_2, \dots, X_k)=0.$$
We will prove a weaker result for the structure $(M_n(\mathbb N), +, \cdot, 0, 1)$ with $n > 1$ instead of $(\mathbb N, +, \cdot, 0, 1)$. It is undecidable whether systems of non-commutative equations have solutions in the set $M_n(\mathbb N)$. Here $0$ is interpreted as the all-zero matrix, while $1$ is interpreted as the unit matrix $I_n$. This matrix satisfies $I_n(a,b) = 1$ iff $a=b$ and $I_n(a,b) = 0$ iff $a \neq b$. Integers $k \in \mathbb N$ arizing as coefficients or as free terms, are always interpreted as $kI_n$. 

An example of diophantine equation solved in $M_2(\mathbb N)$ is shown in \cite{Math1089}. The equation $$AB = 10 \cdot  A + B$$ has the following amazing solution in decimal digits:
$$
\begin{pmatrix}
3 & 4 \\ 8 & 7
\end{pmatrix}
\textcolor{red}
{
\begin{pmatrix}
7 & 2 \\ 4 & 9
\end{pmatrix}
} = 
\begin{pmatrix}
3\textcolor{red}{7} & 4\textcolor{red}{2} \\ 8\textcolor{red}{4} & 7\textcolor{red}{9}
\end{pmatrix}
$$
On the cited internet-site,  other solutions in $M_3(\mathbb N)$ and in $M_4(\mathbb N)$ are shown. In $M_1(\mathbb N)$ the only one non-trivial solution is $A = B = 11$, which unfortunately is not a digit. 

Also, 
$$
{\begin{pmatrix} 1 & 0 \\ 0 & 0 \end{pmatrix}}^k +  {\begin{pmatrix} 0 & 0 \\ 0 & 1 \end{pmatrix}}^k = {\begin{pmatrix} 1 & 0 \\ 0 & 1 \end{pmatrix}}^k 
$$
is a solution of the equation $X^k + Y^k = Z^k$ in $M_2(\mathbb N)$ for all $k \geq 1$. 


Solutions for both equations can be extended to solutions in $M_n(\mathbb N)$ for all $n >2$ by considering $n \times n$ matrices $A'$ and $B'$ whose left upper corners are the matrices $A$ and $B$, respectively $X'$, $Y'$ and $Z'$ whose left upper corners are the matrices $X$, $Y$ and $Z$, all other elements in these matrices being $0$. This is possible only because the equations $AB = 10A + B$, respectively $X^k + Y^k = Z^k$, do not contain any variable-free term. Examples of equations with variable-free terms are shown in \cref{sect:divide} and \cref{sect:undec}.

{\bf Systems of equations} are positive existential propositions of the shape:
$$ S \,\,:\,\,\exists \,X_1, \dots, X_k \,\,\,\,P_1(X_1, \dots, X_k) = 0 \wedge \dots \wedge P_s(X_1, \dots, X_k) = 0$$
where $P_1, \dots, P_s \in \mathbb Z[X_1, X_2, \dots, X_k]$ are non-commutative polynomials.

We denote by $Dph_n(\mathbb N)$ the set of systems which  are satisfiable in $M_n(\mathbb N)$. We define similarly the sets $Dph_n(\mathbb Z)$ and $Dph_n(\mathbb Q)$. By Matiyasevich \cite{Matiyasevich}, the sets $Dph_1(\mathbb N)$ and $Dph_1(\mathbb Z)$ are algorithmically undecidable. The decidability of $Dph_1(\mathbb Q)$ is open, see Koenigsmann \cite{Koenigsmann}.

We prove the following results:

\begin{theorem}\label{diophantinetheories}
Between the sets $Dph_n(\mathbb N)$ and $Dph_m(\mathbb N)$  with $m, n \geq 1$ one has the following relations:
\begin{enumerate}
\item $Dph_n(\mathbb N) = Dph_m(\mathbb N) \longleftrightarrow n = m$.
\item $Dph_n(\mathbb N) \subseteq Dph_m(\mathbb N) \longleftrightarrow n \mid m$.
\end{enumerate}
These are also true for the families  $\left (Dph_n(\mathbb Z)\right )$, respectively $\left (Dph_n( \mathbb Q) \right )$.
\end{theorem}

\begin{theorem}\label{theoremundecidability}
For $n \geq 1$, it is undecidable whether systems of non-commutative polynomial equations:
\begin{eqnarray*}
P_1(X_1, X_2, \dots, X_k) &=&0 \\
P_2(X_1, X_2, \dots, X_k) &=&0 \\
\dots \dots \dots \dots \dots \dots \dots &\dots& \dots \\
P_r(X_1, X_2, \dots, X_k) &=&0
\end{eqnarray*}
have solutions $(X_1, X_2, \dots, X_k) \in M_n(\mathbb N)^k$. 
\end{theorem}



%%%%%%%%%%%%%%%%%%%%%%%%%%%%%%%%%%%%%%%%%%
\section{Solvability depends on matrix-dimension}\label{sect:divide}

For the following classic results, see any classical book touching finitely generated modules over principal ideal commutative rings, like Ion and Radu's Algebra  \cite{IonRadu} or any classic book of linear algebra, like Serge Lang \cite{Lang}. 

Let $K$ be a commutative field. For $f \in K[X]$ and $A \in M_n(K)$, $f = \alpha_0 + \alpha_1X + \dots + \alpha_s X^s$, we define:
$$\xi(f) = f(A) = \alpha_0I_n + \alpha_1 A + \dots + \alpha_s A^s.$$ 
For given matrix $A$, the set $Ker(\xi)$ is an ideal of the principal ideal ring $K[X]$. So there is an element $\mu_A(X) \in K[X]$ such that $Ker(\xi) = (\mu_A)$. The polynomial $\mu_A$ is called {\bf minimal polynomial} of the matrix $A$. 

The polynomial
$$\chi_A(X) := \det(X I_n - A)$$
is called the {\bf characteristic polynomial} of the matrix $A$. By the Cayley-Hamilton Theorem, $\chi_A(A) = 0$. It follows that $\mu_A(X) \mid \chi_A(X)$. Moreover, the polynomials $\mu_A(X)$ and $\chi_A(X)$ have the same irreducible factors in $K[X]$. 

We will also make use of Eisenstein's Irreducibility Criterion. If $P(X) = a_n X^n + a_{n-1} X^{n-1} + \dots + a_0 \in \mathbb Z[X]$ is a polynomial, and $p$ is a prime number such that $p \nmid a_n$, $p \mid a_i$ for $0 \leq i < n$ and $p^2 \nmid a_0$, then $P(X)$ is irreducible in $\mathbb Q[X]$. In particular, the polynomial:
$$ P(X) = X^n - 2 $$
is irreducible in $\mathbb Q[X]$, with $p = 2$. 

\begin{lemma}\label{lemmasolution}
For $n \geq 1$, the following matrix $A \in M_n(\mathbb N)$:
$$ A =
\begin{pmatrix}
0 & 1 & 0 & \dots & 0 \\
0 & 0 & 1 & \dots & 0 \\
\vdots & \vdots & \vdots & \ddots & \vdots \\
0 & 0 & 0 & \dots & 1 \\
2 & 0 & 0 & \dots & 0 
\end{pmatrix}
$$
satisfies the relation $A^n = 2I_n$, so is a solution of the equation $X^n - 2 = 0$ in $M_n(\mathbb N)$. 
\end{lemma} 

\begin{proof}
This works either by direct computation, or by showing that the characteristic polynomial $\chi_A(X) = X^n - 2$. 
\end{proof} 

\begin{lemma}\label{lemmanosmallerdim} 
If $k < n$, then the equation $X^n - 2 = 0$ has no solution in $M_k(\mathbb N)$. 
\end{lemma}

\begin{proof} Let $A \in M_k(\mathbb N)$ be any solution of $X^n - 2$ in this set. The minimal polynomial $\mu_A(X)$ divides $X^n - 2$, which is irreducible, so $\mu_A(X) = X^n-2$. On the other hand, $\mu_A(X)$ divides the characteristic polynomial $\chi_A(X)$. But the degree of the characteristic polynomial is $k < n$, contradiction. 
\end{proof} 

\begin{lemma}\label{lemmadivisibilnecessary}
If $n < m$ and the equation $X^n -2 = 0$ is solvable in $M_m(\mathbb N)$, then $n \mid m$. 
\end{lemma} 

\begin{proof}
Let $A \in M_m(\mathbb N)$ be a solution of $X^n -2 = 0$. Again the minimal polynomial $\mu_A(X)$ divides $X^n-2$ which is irreducible, so $\mu_A(X) = X^n-2$. The characteristic polynomial $\chi_A(X)$ has the same irreducible factors as the minimal polynomial $\mu_A(X)$, so $\chi_A(X) = \mu_A(X)^k$ for some $k > 1$. But the degree of $\chi_A(X)$ is $m$, so $m = kn$. 
\end{proof} 

\begin{lemma}\label{lemmadivisibilsufficient}
If $n \mid m$ then $Dph_n(\mathbb N) \subset Dph_m(\mathbb N)$. 
\end{lemma}

\begin{proof} Let $m = kn$ with $k> 1$. 
Consider the following application $\delta : M_n(\mathbb N) \rightarrow M_m(\mathbb N)$ given by the following $k \times k$ block matrix-representation: 
$$ \delta(A) = 
\begin{pmatrix}
A & O & \dots & O \\
O & A & \dots & O \\
\vdots & \vdots & \ddots & \vdots \\
O & O & \dots & A
\end{pmatrix}
$$
where $O \in M_n(\mathbb N)$ is the zero matrix. We observe that $\delta$ is an injective homomorphism according to the language $(+, \cdot, 0, 1)$. 

If $(x_1, \dots, x_k) \in M_n(\mathbb N)^k$ is a solution of the system of equations:
$$P_1(x_1, x_2, \dots, x_k) = 0 \wedge \dots \wedge P_s(x_1, \dots, x_k) = 0, $$
then $(\delta(x_1), \delta(x_2), \dots, \delta(x_n)) \in M_m(\mathbb N)^k$ is a solution of the same system of equations in $M_m(\mathbb N)$. 
\end{proof} 

Now we are ready to prove \cref{diophantinetheories}. For (1) if $n = m$ then $Dph_n (\mathbb N) = Dph_m(\mathbb N)$. Also, if $n \neq m$, say $n < m$, then the equation $X^m - 2 = 0$ belongs to $Dph_m(\mathbb N)$ by \cref{lemmasolution} but does not belong to $Dph_n (\mathbb N)$ by \cref{lemmanosmallerdim}, so these sets are different. For (2), suppose that $Dph_n(\mathbb N) \subseteq Dph_m(\mathbb N)$. Again by \cref{lemmasolution}, the equation $X^n -2 = 0$ belongs to $Dph_n(\mathbb N)$. As it belongs also to $Dph_m(\mathbb N)$, by the \cref{lemmadivisibilnecessary} one gets that $n \mid m$. Also, if $n \mid m$, by \cref{lemmadivisibilsufficient} we get $Dph_n(\mathbb N) \subseteq Dph_m(\mathbb N)$.

For the  families  $\left (Dph_n(\mathbb Z)\right )$, respectively $\left (Dph_n( \mathbb Q) \right )$, the proof works similarly.  \qed




%%%%%%%%%%%%%%%%%%%%%%%%%%%%%%%%%%%%%%%%%%
\section{Solvability is undecidable}\label{sect:undec}

 For $1 \leq i, j \leq n$, the {\bf elementary matrix} $E_{i,j}$ is defined by $E_{i,j}(a,b) = 1$ iff $i = a$ and $j = b$, respectively $E_{i,j}(a,b) = 0$ otherwise. The matrices $E_{ii}$ are called {\bf elementary diagonal} matrices. 

\begin{lemma}\label{lemmadefinitionelementarydiagonal}
Let $(Y, A_1, A_2, \dots, A_{n-1})$ be any solution of the system:
   \begin{eqnarray*}
       Y + A_1YA_1 + A_2YA_2 + \dots + A_{n-1} Y A_{n-1} & = & 1 \\
       A_1^2 A_2^2 \dots A_{n-1}^2 & = & 1 
   \end{eqnarray*}
in $M_n(\mathbb N)$. Then there is some $1 \leq i \leq n$ such that $Y = E_{i,i}$. Moreover, for every choice $Y = E_{i,i}$ there is a solution $(Y, A_1, \dots, A_{n-1})$ satisfying this choice.
\end{lemma} 

\begin{proof}
For $n=1$ there is no $A_i$ and the system has the solution $Y = 1 = E_{1,1}$. Take $n \geq 2$. We observe that all matrices $A_i$ are invertible. If any product $A_k YA_k = 0$, then $Y = 0$ and the whole left-hand side of the first equation would be $0$. So all the matrices $Y$, $A_1YA_1$, $\dots$, $A_{n-1}YA_{n-1}$ are different from $0$. On the other hand, $M_n(\mathbb N)$ is closed for multiplication, so all these $n$ many matrices consist of natural numbers only. As $1$ is interpreted as the unit matrix $I_n$, the only one way one can decompose $I_n$ in $M_n(\mathbb N)$ as a sum of $n$ many non-zero matrices is:
$$I_n = E_{1,1} + E_{2,2} + \dots + E_{n,n}.$$
It follows that the following sets are identic:
$$\{Y, A_1YA_1, \dots, A_{n-1}YA_{n-1}\} = \{E_{1,1}, E_{2,2}, \dots, E_{n,n}\}.$$
This shows that in any solution, there is some $1 \leq i \leq n$ such that $Y = E_{i,i}$. Now we have to show that a solution always exists. Well, for every possible choice $Y = E_{i,i}$ many solutions do exist. For example, if we want that $A_1YA_1 = E_{j,j}$ for some $1 \leq j \leq n$ with $j \neq i$, we choose $A_1$ to be the permutation matrix corresponding to the transposition $(i,j)$. The action of this matrix on the canonical orthogonal basis $e_1 = (1, 0, \dots, 0)$, $\dots$, $e_n = (0, 0, \dots, 1)$ is the following:  $A_1 e_i = e_j$, $A_1e_j = e_i$ and $A_1 e_k = e_k$ for $k \notin \{i,j\}$. Indeed, we see that:
$$A_1 E_{i,i} A_1 = E_{j,j},$$
and also $A_1^2 = 1$. So all $n-1$ many matrices $A_1$, $\dots$, $A_{n-1}$ are chosen to be the different permutation matrices corresponding to the transpositions $(i,j)$ with $j \neq i$, taken in some arbitrary order. All these matrices satisfy $A_k^2 = 1$ for $1 \leq k \leq n-1$, so the condition $A_1^2A_2^2 \dots A_{n-1}^2 = 1$ is fulfilled. 
\end{proof} 

For example, if $n = 2$, a possible solution is:
$$Y= \begin{pmatrix} 0 & 0 \\ 0 & 1 \end{pmatrix}= E_{2,2},\,\,\,\, A_1 = \begin{pmatrix} 0 & 1 \\ 1 & 0 \end{pmatrix} = (1,2).$$

For some fixed $1 \leq i \leq n$, consider the set:
$$\Sigma_i = \{ A \in M_n(\mathbb N) \,|\, E_{i,i} A = A E_{i,i} \}.$$ 
We call $\pi_{i,i}: M_n(\mathbb N) \rightarrow \mathbb N$ the application given by $\pi_{i,i}(A) = A(i,i)$. 

\begin{lemma}\label{lemmasigmai}
The set $\Sigma_i$ can be characterized as follows:
$$\Sigma_i = \{ A \in M_n(\mathbb N) \,|\,1 \leq k \leq n \,\wedge \, k \neq i \rightarrow A(k,i) = A(i,k) = 0 \}.$$ 
So $A(i,i)$ is the only one element on row $i$ and on column $i$, which is not necessarily $0$. 
\end{lemma}

\begin{proof} If the matrix $A$ consists of the elements $A(k,j)$, one has:
$$ A E_{i,i} = 
\begin{pmatrix}
A(1,1) & A(1,2) & \dots & A(1,n) \\
\vdots & \vdots & \vdots & \vdots \\
A(i,1) & A(i,2) & \dots & A(i,n) \\
\vdots & \vdots & \vdots & \vdots \\
A(n,1) & A(n,2) & \dots & A(n,n) 
\end{pmatrix}
\begin{pmatrix}
0 & \dots & 0 & \dots & 0\\
\vdots & \vdots & \vdots & \vdots & \vdots \\
0 & \dots &1 & \dots & 0\\
\vdots & \vdots & \vdots & \vdots & \vdots \\
0 & \dots & 0 & \dots & 0
\end{pmatrix} = 
\begin{pmatrix}
0 & \dots & A(1,i) & \dots & 0\\
\vdots & \vdots & \vdots & \vdots & \vdots\\
0 & \dots &A(i,i) & \dots & 0\\
\vdots & \vdots & \vdots & \vdots & \vdots \\
0 & \dots & A(n,i) & \dots & 0
\end{pmatrix}
$$

$$ E_{i,i} A = 
\begin{pmatrix}
0 & \dots & 0 & \dots & 0\\
\vdots & \vdots & \vdots & \vdots & \vdots \\
0 & \dots &1 & \dots & 0\\
\vdots & \vdots & \vdots & \vdots & \vdots \\
0 & \dots & 0 & \dots & 0
\end{pmatrix}
\begin{pmatrix}
A(1,1) & A(1,2) & \dots & A(1,n) \\
\vdots & \vdots & \vdots & \vdots \\
A(i,1) & A(i,2) & \dots & A(i,n) \\
\vdots & \vdots & \vdots & \vdots \\
A(n,1) & A(n,2) & \dots & A(n,n) 
\end{pmatrix}
 = 
$$
$$ =
\begin{pmatrix}
0 & \dots & 0 & \dots & 0\\
\vdots & \vdots & \vdots & \vdots & \vdots\\
A(i,1) & \dots &A(i,i) & \dots & A(i,n)\\
\vdots & \vdots & \vdots & \vdots & \vdots \\
0 & \dots & 0 & \dots & 0
\end{pmatrix}
$$ 
From the equality $AE_{i.i} = E_{i,i} A$, the conclusion follows.
\end{proof}
The typical element of $\Sigma_i$ has the following shape:
$$
\begin{pmatrix}
A(1,1) & A(1,2) & \dots &A(1, i-1) & 0 & A(1, i+1) & \dots & A(1,n) \\
\vdots & \vdots & \vdots & \vdots & \vdots & \vdots &\vdots & \vdots \\
A(i-1,1) & A(i-1,2) & \dots &A(i-1, i-1) & 0 & A(i-1, i+1) & \dots & A(i-1,n) \\
&&&&&&& \\
0 & 0 & \dots &0 & A(i,i) & 0 & \dots & 0  \\
&&&&&&& \\
A(i+1,1) & A(i+1,2) & \dots &A(i+1, i-1) & 0 & A(i+1, i+1) & \dots & A(i+1,n) \\
\vdots & \vdots & \vdots & \vdots & \vdots & \vdots &\vdots & \vdots  \\
A(n,1) & A(n,2) & \dots & A(n,i-1) & 0  & A(n, i+1) & \dots & A(n,n)
\end{pmatrix}
$$
Here, the values denoted by $A(j,k)$ with various indexes, might be also equal $0$. But only the values marked by $0$ must be equal $0$ in every element of $\Sigma_i$. 

\begin{lemma}\label{lemmahomomorphism} For all $1 \leq i \leq n$, one has:
\begin{enumerate}
\item $(\Sigma_i, +, \cdot, 0, I_n)$ is a substructure of $(M_n(\mathbb N), +, \cdot, 0, I_n)$.  
\item $\pi_{i,i} : (\Sigma_i, +, \cdot, 0, I_n) \rightarrow (\mathbb N, +, \cdot, 0, 1)$ is a surjective homomorphism. 
\item $\alpha :  (\mathbb N, +, \cdot, 0, 1) \rightarrow (\Sigma_i, +, \cdot, 0, I_n) $ given by $\alpha(a) = aI_n$ is an injective homomorphism.
\item $\pi_{i,i} \circ \alpha = id_{\mathbb N}$.
\end{enumerate} 
\end{lemma}

\begin{proof} By direct computation, using the characterisation of $\Sigma_i$ given in \cref{lemmasigmai}. \end{proof} 

\begin{lemma}\label{lemmaembedding}
For $k \geq 1$ and $f \in \mathbb Z[x_1, \dots, x_k]$, the equation $f(x_1, x_2, \dots, x_k) = 0$ has a solution  $(x_1, x_2, \dots, x_k) \in \mathbb N^k$ if and only if the following system of $k+3$ equations in $n+k$ unknowns: 

  \begin{eqnarray*}
       Y + A_1YA_1 + A_2YA_2 + \dots + A_{n-1} Y A_{n-1} & = & 1 \\
       A_1^2 A_2^2 \dots A_{n-1}^2 & = & 1 \\
       X_1 Y - Y X_1 &=& 0\\
       X_2 Y - Y X_2 &=& 0\\
      \dots \dots  \dots \dots \dots \dots \dots & \dots & \dots \\
       X_k Y -  Y X_k &=& 0\\
      f(X_1, X_2, \dots, X_k) &= & 0
   \end{eqnarray*} 

has a solution $(Y, A_1, A_2, \dots, A_{n-1}, X_1, X_2, \dots, X_k) \in M_n(\mathbb N)^{n+k}$. 
\end{lemma}

\begin{proof}
Let $(x_1, x_2, \dots, x_k) \in \mathbb N^k$ be a solution of the equation $f(x_1, x_2, \dots, x_k) = 0$. Then the matrices $X_j = x_j I_n$ build a solution of $f(X_1, X_2, \dots, X_k) = 0$ in $M_n(\mathbb N)$. We know that the system built up by the first two equations is always solvable, by \cref{lemmadefinitionelementarydiagonal}. For every solution $(Y, A_1, \dots, A_{n-1})$ of the system consisting of the first two equations, one has $X_j Y = Y X_j$ because the diagonal matrices (also called scalar matrices) commute with every matrix. 

Now let $(Y, A_1, A_2, \dots, A_{n-1}, X_1, X_2, \dots, X_k) \in M_n(\mathbb N)^{n+k}$ be a solution of the system. Again by  \cref{lemmadefinitionelementarydiagonal}, there is an $1 \leq i \leq n$ such that $Y = E_{i,i}$. Because of the equations $X_j Y = YX_j$ for all $1 \leq j \leq k$, we find that $X_1, X_2, \dots X_k \in \Sigma_i$. Using \cref{lemmahomomorphism}, we observe that the tuple $(\pi_{i,i}(X_1), \pi_{i,i}(X_2), \dots, \pi_{i,i}(X_k)) \in \mathbb N^k$ is a solution for the equation $f(x_1, x_2, \dots, x_k) = 0$. 
\end{proof}

Now we can prove \cref{theoremundecidability}. For $n = 1$, we get the Theorem of Matiyasevich, see \cite{Matiyasevich}. For $n > 1$, we apply the Theorem of Matiyasevich and \cref{lemmaembedding}. \qed

%%%%%%%%%%%%%%%%%%%%%%%%%%%%%%%%%%
\section{Other results}\label{sect:other}

The first category of results are corollaries to \cref{theoremundecidability}. A subset $B \subset \mathbb N$ is called {\bf additive basis} of the set of real numbers if there is an $d = d_B$ such that for every $x \in \mathbb N$ there are $x_1, x_2, \dots, x_d \in B$ such that 
$$x = x_1 + x_2 + \dots + x_d.$$

\begin{corollary}
Let $B \subset \mathbb N$ be an additive basis for the set of natural numbers, such that $0, 1 \in B$. For every $n \geq 1$, it is undecidable whether systems of non-commutative polynomials have solutions in the set $M_n(B)$. 
\end{corollary}

\begin{proof}
It follows that for every matrix $C \in M_n(\mathbb N)$ there are matrices $C_1$, $\dots$, $C_d$ $\in M_n(B)$ such that:
$$ C = C_1 + C_2 + \dots + C_d.$$
One takes a system $S$ of non-commutative polynomials and built a system $S'$ by replacing every variable $X$ with a sum ${^1X} + {^2X} + \dots + {^dX}$, where $^1X$, $^2X$, $\dots$, $^dX$ are new variables, used such that their sum replaces only the variable $X$ and no other variables. $S$ is solvable in $M_n(\mathbb N)$ iff $S'$ is solvable in $M_n(B)$. It is important that $0, 1 \in B$ because the unknown $Y$ can be satisfied only with matrices containing zeros and a one.\end{proof} 

For example, for all $t \geq 2$, the set of $t$-powers is an additive basis of $\mathbb N$ (see Hilbert \cite{Hilbert} for the original proof, or Khinchin \cite{Khinchin} for a modern proof), so for every $t \geq 2$ it is undecidable whether polynomial systems of equations have solutions consisting of matrices filled with $t$-powers. The same is true for polygonal numbers (see Nathanson's paper \cite{Nathanson} or his book \cite{Nathanson1}). As the set of prime numbers completed with $0$ and $1$ also builts an additive basis for $\mathbb N$ (see Schnirelmann \cite{Schnirelmann}, Vinogradov \cite{Vinogradov} or Nathanson's book \cite{Nathanson1}), it is undecidable whether solutions consisting of matrices filled with $0$, $1$ and prime numbers do exist. Such deductions have been already done by the author in \cite{Prunescu} for the case of equations over $\mathbb N$.

If one could find a positive existential definition of $M_n(\mathbb N)$ in $M_n(\mathbb Z)$, one could transfer \cref{theoremundecidability} to the rings $M_n(\mathbb Z)$. Also, such a result would be proved if one finds a positive existential definition of the set $\{E_{1,1}, E_{2,2}, \dots, E_{n,n}\}$ in $M_n(\mathbb Z)$. We also do not know how to give diophantine definitions for the conjunction and the disjunction of equations, in order to transfer \cref{theoremundecidability} from systems of equations to equations. 

It is much easier to prove undecidability if we allow equations with parameters in $M_n(\mathbb N)$. In this case the result works directly for equations. We do not need considering systems of equations anymore. 

\begin{theorem}
For $n \geq 1$, it is undecidable, whether non-commutative equations with parameters in $M_n(\mathbb Z)$ have solutions in $M_n(\mathbb N)$. It is also undecidable, whether  non-commutative equations with parameters in $M_n(\mathbb Z)$ have solutions in $M_n(\mathbb Z)$.  Moreover, it is sufficient to consider parameters in $M_n(\{0,1\})$. 
\end{theorem} 

\begin{proof}
We observe that for every matrix $A\in M_n(\mathbb Z)$, one has:
$$
E_{1,1} A E_{1,1} =
\begin{pmatrix}
A(1,1) & 0 & \dots & 0 \\
0 & 0 & \dots & 0 \\
\vdots & \vdots & \ddots & \vdots \\
0 & 0 & \dots & 0 
\end{pmatrix} = A(1,1)E_{1,1}.
$$
These matrices build a ring with $1$, which is a subring of the ring without $1$ given by $(M_n(\mathbb Z), +, \cdot, 0)$. For every polynomial $f \in \mathbb Z[x_1, x_2, \dots, x_k]$ we construct a polynomial $\tilde f \in \mathbb Z[X_1, X_2, \dots, X_k, E]$ as follows. To every monomial $x_{i_1}x_{i_2} \dots x_{i_s}$ of $f$ we construct the new monomial $$EX_{i_1}EX_{i_2}E\cdots  EX_{i_s}E.$$ Between every two variables, a copy of the new non-commutative variable $E$ is inserted. Also, copies of $E$ are written as prefix and as suffix of the monomial. The intention is to interpret $E$ as the parameter $E_{1,1} \in M_n(\{0,1\})$ and to apply the fact that $E_{1,1}^2 = E_{1,1}$. 
Suppose that we denote $X_i(1,1)$ as new numeric variable $x_i$. It follows that: 
$$EX_{i_1}EX_{i_2}E\cdots EX_{i_s}E = (EX_{i_1}E)(EX_{i_2}E) \cdots (EX_{i_s}E) = x_{i_1}x_{i_2} \cdots x_{i_s}E.$$
Also, the free term $kI_n$ is replaced in $\tilde f$ by $kE$. 

We observe that the equation $f(x_1, x_2, \dots, x_k) = 0$ has a solution in $\mathbb N$ if and only if the equation $\tilde f(X_1, X_2, \dots, X_k, E)=0$ has a solution in $M_n(\mathbb N)$ with $E = E_{1,1}$. Also,  $f(x_1, x_2, \dots, x_k) = 0$ has a solution in $\mathbb Z$ if and only if the equation $\tilde f(X_1, X_2, \dots, X_k, E)=0$ has a solution in $M_n(\mathbb Z)$ with $E = E_{1,1}$. 

Indeed, consider $\tilde f$ as polynomial in variables $X_1, X_2, \dots, X_k$ with coefficients and parameters in $M_n(\mathbb Z)$. If the tuple $(a_1, a_2, \dots a_k) \in \mathbb N^k$ is a solution of the equation $f(x_1, x_2, \dots, x_k) = 0$, then $(a_1E_{1,1}, a_2E_{1,1}, \dots, a_kE_{1,1})$ is a solution of the equation $\tilde f(X_1, X_2, \dots,X_k) = 0$ in $M_n(\mathbb N)$. On the other hand, if the tuple $(A_1, A_2, \dots, A_k) \in M_n(\mathbb N)^k$ satisfies $\tilde f(X_1, X_2, \dots, X_n)=0$, then the tuple $(\pi_{1,1}(A_1), \pi_{1,1}(A_2), \dots, \pi_{1,1}(A_k)) \in \mathbb N^k$  satisfies $f(x_1, x_2, \dots, x_k) = 0$. All remains true if we replace $\mathbb N$ with $\mathbb Z$. 
\end{proof}

It is known that any matrix $X \in M_n(\mathbb Z)$ is equal to $aI_n$ for some $a \in \mathbb Z$ if\footnote{Indeed, if $X$ commutes with all $E_{i,i}$, then it is a diagonal matrix. If, moreover, it commutes with the all-one matrix $\sum E_{i,j}$, then it is a scalar matrix.}  and only if
$$\forall \,A\,\,\,\, XA = AX.$$
 If we combine this information with the fact that Hilbert's Tenth Problem has a negative answer for the ring $\mathbb Z$, we conclude that the $\exists \forall$-theory of $M_n(\mathbb Z)$ is undecidable. 

Non-commutative equations $P(x_1, x_2, \dots, x_k) = 0$ with $P \in \mathbb Z[X_1, X_2, \dots, X_k]$ can be always written in the equivalent form:
$$ P_1(x_1, x_2, \dots, x_k) = P_2(x_1, x_2, \dots, x_k), $$
where on both sides one has only additions and multiplications. This is so, because negative monomials from the left-hand side are moved with changed sign to the right-hand side. 

\begin{proposition}\label{prop}
Let $(\mathcal A, +, \cdot, 0, 1)$ and $(\mathcal B, +, \cdot, 0, 1)$ be two infinite structures with the same signature. Suppose that there exist two homomorphisms $\alpha: \mathcal A \rightarrow \mathcal B$ and $\beta : \mathcal B \rightarrow \mathcal A$ according to this signature. Then the following are equivalent: (1) it is decidable whether  equations from $\mathbb Z[x_1, x_2, \dots, x_k]$ have solutions in $\mathcal A$ and (2) it is decidable whether  equations from $\mathbb Z[x_1, x_2, \dots, x_k]$ have solutions in $\mathcal B$.
\end{proposition} 

\begin{proof}
An equation is solvable in $\mathcal A$ if and only if it is solvable in $\mathcal B$. Indeed, in the first case, if $(x_1, x_2, \dots, x_k) \in \mathcal A^n$ is a solution of an equation in $\mathcal A$, then $(\alpha(x_1), \alpha(x_2), \dots, \alpha(x_k))\in \mathcal B^n$ is a solution of the same equation  in $\mathcal B$. From $\mathcal B$ to $\mathcal A$ the transfer is done with $\beta$. 
\end{proof}

Using \cref{prop} we can prove that equation solvability is undecidable for many subrings of $M_n(\mathbb Z)$. In the next examples we take always $\mathcal A = \mathbb Z$ and $\alpha(k) = kI_n$. The structures $\mathcal B$ and the homomorphisms $\beta$ can be chosen as follows:

\begin{enumerate}
\item $\mathcal B = D_n(\mathbb Z)$, the diagonal matrices with integer coefficients, where the elements on the main diagonal are free to take any value, but all other elements are $0$. Any projection $\pi_{k,k}$ is good to play the role of $\beta$.
\item $\mathcal B = T_n(\mathbb Z)$, the upper triangular matrices with integer coefficients. They are defined by the condition that for all $1 \leq i, j \leq n$, if $i > j$, then $A(i, j)=0$. Again any projection  $\pi_{k,k}$ is good to play the role of $\beta$. 
\item  $\mathcal B = \Sigma_i$ defined by: $1 \leq k \leq n \,\wedge \, k \neq i \rightarrow A(k,i) = A(i,k) = 0$. Only $A(i,i)$ is allowed to be non-zero on both the $i$-th row and the $i$-th column; $\beta = \pi_{i,i}$. 
\item  $\mathcal B = \Gamma_i$ defined by: $1 \leq k \leq n \,\wedge \, k \neq i \rightarrow A(k,i) = 0$. Only $A(i,i)$ is allowed to be non-zero on the $i$-th column; $\beta = \pi_{i,i}$. 
\item  $\mathcal B = \Lambda_i$ defined by: $1 \leq k \leq n \,\wedge \, k \neq i \rightarrow A(i,k) = 0$. Only $A(i,i)$ is allowed to be non-zero on the $i$-th row; $\beta = \pi_{i,i}$.
\item $\mathcal B = R_i$ defined by: $$1 \leq k,j \leq n\, \wedge \, k \geq i \, \wedge \, j \leq i \,\wedge \, (k,j) \neq (i,i) \rightarrow A(k,j) = 0.$$
Only $A(i,i)$ is allowed to be non-zero in the rectangle $[i,n] \times [1,i]$; $\beta = \pi_{i,i}$. 
\item $\mathcal B = RR_i$ defined by: 
 $$1 \leq k,j \leq n\, \wedge \, [ (k \geq i \, \wedge \, j \leq i) \, \vee  \,(k \leq i \, \wedge \, j \geq i)\,]\, \wedge \, (k,j) \neq (i,i) \rightarrow A(k,j) = 0.$$
Only $A(i,i)$ is allowed to be non-zero in the union of two rectangles $([i,n] \times [1,i]) \cup ([1,i] \times [i,n])$; $\beta = \pi_{i,i}$.
\end{enumerate}

%%%%%%%%%%%%%%%%%%%%%%%%%%%%%%%%%%%%
\section{On equations without free term}\label{sect:nofreeterm}

One can define the following problems:

$$ H_n(\mathbb Z) = \{ f \in \mathbb Z[X_1, X_2, \dots, X_k]\,|\,f(\lambda X_1, \lambda X_2, \dots, \lambda X_k) = \lambda^{\deg(f)} f(X_1, X_2, \dots, X_k) \,\wedge  $$ $$ \wedge \, \left (\exists \,(Y_1, Y_2, \dots, Y_k) \in M_n(\mathbb Z)^k\,\,\,\,(Y_1, Y_2, \dots, Y_k) \neq (0, 0, \dots, 0) \,\wedge \, f(Y_1, Y_2, \dots, Y_k) = 0 \right )\}, $$ 

$$ F_n(\mathbb Z) = \{ f \in \mathbb Z[X_1, X_2, \dots, X_k]\,|\,f(0, 0, \dots, 0) = 0 \,\wedge  $$ $$ \wedge \, \left (\exists \,(Y_1, Y_2, \dots, Y_k) \in M_n(\mathbb Z)^k\,\,\,\,(Y_1, Y_2, \dots, Y_k) \neq (0, 0, \dots, 0) \,\wedge \, f(Y_1, Y_2, \dots, Y_k) = 0 \right )\}, $$ 

The first problem is about finding non-trivial solutions for {\bf homogeneous}  equations. The second one is about finding non-trivial solutions for equations {\bf without free term}. Both problems can be equally defined for $\mathbb N$ and $\mathbb Q$ instead of $\mathbb Z$. Also, all these problems can be put in the version $SH_n(\mathbb Z)$ and $SF_n(\mathbb Z)$ for systems of homogeneous equations respectively for systems of equations without free terms, and this can be done as well for  $\mathbb N$ and $\mathbb Q$ instead of $\mathbb Z$. 

While the decidability of $H_1(\mathbb Z)$ is equivalent with the decidability of $Dph_1(\mathbb Q)$ (see Matiyasevich \cite{Matiyasevich}) and is so far an open problem, for all other problems defined here even less information is known. However, an observation can be done:

\begin{proposition}
Let $( G_n )$ be one of the $12$ problems defined above. Then the sequence $( G_n )$ is monotone:
$$ n \leq m \longleftrightarrow G_n \subseteq G_m .$$
\end{proposition}

\begin{proof}
Keep in mind the problem $F_n(\mathbb Z)$. For all the other problems mentioned, the proof is similar. For $n \leq m$, we consider the application:
$$\gamma : M_n(\mathbb Z) \rightarrow M_m(\mathbb Z), \,\,\,\,\,\, \gamma(A) = \begin{pmatrix} A & ^T O_1 \\ O_1 & O_2 \end{pmatrix}.$$
In this block representation, $O_1$ is the $(m-n) \times n$ matrix consisting of zeros, $^T O_1$ is the $n \times (m-n)$ matrix consisting of zeros, and $O_2$ is the $(m-n) \times (m-n)$ matrix consisting of zeros. We observe that $\gamma$ is an injective homomorphism of rings without $1$. If $f(0, 0, \dots, 0) = 0$, and $(Y_1, Y_2, \dots, Y_k) \in M_n(\mathbb Z)^k$, $(Y_1, Y_2, \dots, Y_k) \neq (0,0,\dots , 0)$ and $f(Y_1, Y_2, \dots, Y_k) =0$, then $(\gamma(Y_1), \gamma( Y_2), \dots, \gamma(Y_k) )  \in M_m(\mathbb Z)^k$, $(\gamma(Y_1), \gamma( Y_2), \dots, \gamma(Y_k) )  \neq (0,0,\dots , 0)$ and $f(\gamma(Y_1), \gamma( Y_2), \dots, \gamma(Y_k) )  = 0$. 
\end{proof}

\begin{thebibliography}{99} 

\bibitem{Cauchy}
{\bf A. L. Cauchy}
{\it Demonstration du theoreme general de Fermat sur les nombres polygones.}
 Mem. Sci. Math. Phys. Inst.  France 14, 1, 177 – 220, 1813.

\bibitem{Demeyer}
{\bf J. Demeyer}
{\it Problems related to extensions of Hilbert's Tenth Problem.}
Moderated by B. Poonen and T. Scanlon, notes by J. Demeyer,

https://aimath.org/WWN/hilberts10th/hilberts10th.pdf 

\bibitem{Hilbert}
{\bf David Hilbert}
{\it Beweis für die Darstellbarkeit der ganzen Zahlen durch eine feste Anzahl n-ter Potenzen (Waringsches Problem).}
 Mathematische Annalen, 67, 3,  281 -- 300, 1909.

\bibitem{IonRadu}
{\bf Ion D. Ion, Nicolae Radu}
{\it Algebra.}
Editura Didactica si Pedagogica, Bucuresti, 1981.

\bibitem{Khinchin}
{\bf Aleksandr Yakovlevich Khinchin}
{\it Three pearls in number theory.}
Dover Publications, Mineola, New York, 1998.

\bibitem{Koenigsmann}
{\bf Jochen Koenigsmann}
{\it Defining $\mathbb Z$ in $\mathbb Q$.}
Annals of Mathematics 183, 73 -- 93, 2016. 


\bibitem{Lang}
{\bf Serge Lang}
{\it Linear Algebra (3rd ed.).}
Springer, 1987.

\bibitem{Matiyasevich}
{\bf Yuri V. Matiyasevich}
{\it Hilbert's tenth problem.}
 MIT Press, 1993. 

\bibitem{Math1089}
{\bf Math1089}
{\it Matrices are beautiful.}
https://math1089.in/matrices-are-beautiful/ 


\bibitem{Nathanson}
{\bf Melvyn B. Nathanson}
{\it A Short Proof of Cauchy's Polygonal Number Theorem.}
Proceedings of the American Mathematical Society, 
99, 1,  22 -- 24, 1987. 

\bibitem{Nathanson1}
{\bf Melvyn B. Nathanson}
{\it Additive Number Theory -- The Classical Bases.}
Graduate Texts in Mathematics 164,
Springer Verlag, New York, 1997. 

\bibitem{Prunescu}
{\bf Mihai Prunescu}
{\it Undecidable and decidable restrictions of Hilbert's Tenth Problem: images of polynomials vs. images of exponential functions.}
Mathematical Logic Quarterly, 52, 1, 14 -- 19, 2006.

\bibitem{Schnirelmann}
{\bf Lev Genrikhovich Schnirelmann:}
{\it  \"Uber additive Eigenschaften von Zahlen.}
Mathematische Annalen 107, 649 – 690, 1933. 

\bibitem{Vinogradov}
{\bf Ivan Matveevich Vinogradov}
{\it Representation of an odd number as the sum of three primes.} 
Doklady Akad. Nauk SSSR 15, 6 -- 7, 291 -- 294,  1937. 
English translation in: Selected Works, 101 -- 106, Springer Verlag, Berlin, 1985.
 

\end{thebibliography}


\end{document} 
